% Figure: energy-flow schematics for a heat engine and a refrigerator/heat
% pump, drawn side by side. Hot reservoir on top, cold reservoir on bottom,
% the device in the middle with arrows showing Q_h, Q_c, and W.
\begin{figure}[htbp]
\centering
\begin{tikzpicture}[
  resv/.style={draw, very thick, minimum width=3.2cm, minimum height=0.8cm, font=\small},
  dev/.style={draw, very thick, fill=englime!15, minimum width=2.0cm, minimum height=1.2cm, font=\small},
  arr/.style={-{Latex[length=2.5mm]}, very thick},
]
% Heat engine (left)
\node[resv, fill=engred!12] (he) at (0,3) {hot reservoir $T_h$};
\node[resv, fill=engblue!12] (ce) at (0,-3) {cold reservoir $T_c$};
\node[dev] (eng) at (0,0) {engine};
\draw[arr, engred] (he.south) -- node[right, font=\scriptsize] {$Q_h$} (eng.north);
\draw[arr, engblue] (eng.south) -- node[right, font=\scriptsize] {$Q_c$} (ce.north);
\draw[arr, engorange] (eng.east) -- node[above, font=\scriptsize] {$W$} ++(1.8,0);
\node[font=\scriptsize, anchor=north] at (0,-3.9) {(a) heat engine: $W = Q_h - Q_c$};
% Refrigerator / heat pump (right)
\begin{scope}[xshift=7cm]
\node[resv, fill=engred!12] (hr) at (0,3) {hot reservoir $T_h$};
\node[resv, fill=engblue!12] (cr) at (0,-3) {cold reservoir $T_c$};
\node[dev] (ref) at (0,0) {refrigerator};
\draw[arr, engred] (ref.north) -- node[right, font=\scriptsize] {$Q_h$} (hr.south);
\draw[arr, engblue] (cr.north) -- node[right, font=\scriptsize] {$Q_c$} (ref.south);
\draw[arr, engorange] (ref.east) ++(1.8,0) -- node[above, font=\scriptsize] {$W$} (ref.east);
\node[font=\scriptsize, anchor=north] at (0,-3.9) {(b) refrigerator: $Q_h = Q_c + W$};
\end{scope}
\end{tikzpicture}
\caption[Energy-flow schematics: heat engine and refrigerator]{Energy-flow
diagrams. (a) A heat engine draws $Q_h$ from the hot reservoir, delivers work
$W$, and rejects $Q_c$ to the cold reservoir, with $W = Q_h - Q_c$ by the first
law. (b) A refrigerator or heat pump runs the arrows in reverse: input work $W$
drives heat $Q_c$ out of the cold reservoir and dumps $Q_h = Q_c + W$ into the
hot one. The cooling coefficient of performance is $Q_c/W$; the heating
coefficient of performance is $Q_h/W$, larger by exactly one.}
\label{fig:vol04ch03:engine-refrigerator-flow}
\end{figure}
